\chapter{Conclusion}
\label{ch:conclusion}

This thesis developed a hybrid detector for sensor attacks under incomplete observations and evaluated it on two water networks. The experiments brought out a recurring problem: an expert can score well on its own and still add little to the final detector. Its value depends on the evidence it adds and on how its alarms are combined with the rest of the system. Redundant detections do not compensate for extra false positives.

The final design separates General, Drift and Noise. General combines five residual classifiers using a target-group-excluding reference, while Drift and Noise use temporal evidence with a declared three-hour delay. Routing, feedback and verification check the specialist alarms and preserve General's immediate path. This arrangement developed through graph and temporal studies, expert diagnostics and candidates that were tried and rejected.

In Modena, guarded integration raised pooled F1 from 0.6863 to 0.8672 and reduced false positives by 84.2\%, at the cost of small average drift and noise losses. L-Town benefited from a different change. Adding received-pressure history inside General improved every family mean, pooled F1 and mean clean-period FPR in the final matched comparison. Both the representation and the acceptance of alarms affected full-system performance.

The final evaluation used ten model seeds and six sources per network. Modena had mean pooled F1 of 0.8669 and family macro F1 of 0.9045; all five family means exceeded 0.80. L-Town reached 0.8950 and 0.9083, respectively, with four family means above 0.89 and replay at 0.7312. Clean-period false-positive rates were 0.0584\% and 0.0550\% for Modena and L-Town. Additional fits preserved these broad profiles.

Similar macro scores should not obscure the remaining weaknesses. Modena's drift and noise results vary by source. In L-Town, history improves every paired replay result but does not reach the later 0.75 objective. The larger network supports strong detection of the other families while leaving this replay problem unresolved.

The single-classifier controls put the added complexity in context under the same maximum three-hour allowance. The hybrid improves pooled F1 and temporal-family coverage in Modena. In L-Town, the single classifier has slightly better pooled F1 and fewer clean false alarms, but the hybrid detects replay, drift and noise better in every paired cell. The preferable design depends on the family coverage required. When background precision has priority, the simpler control remains a credible option.

The reported results apply to generated, observed pressure endpoints under fixed attack and missingness conditions. Fits share evaluation sources, and specialist decisions keep their stated delay. Within those conditions, the detector combines broad family coverage with controlled background errors, with the remaining costs reported alongside the gains. Operational use still requires incident-level evaluation, wider observation conditions and measurements of computational cost.
